\chapter{PENDAHULUAN}

\section{Latar Belakang}

Dalam era digital yang semakin berkembang, pengelolaan infrastruktur IT mengalami transformasi signifikan menuju pendekatan yang lebih otomatis dan terstruktur. \textit{Kubernetes} telah menjadi standar \textit{de facto} untuk orkestrasi kontainer, memungkinkan organisasi untuk men-deploy, mengelola, dan menskalakan aplikasi dengan efisien. Namun, kompleksitas pengelolaan klaster Kubernetes mendorong lahirnya praktik \textit{GitOps} yang mengubah cara tim operasional dan pengembangan berkolaborasi \citep{weaveworks-gitops}.

GitOps adalah paradigma operasional yang menggunakan repositori Git sebagai sumber kebenaran tunggal untuk infrastruktur dan aplikasi. Dengan GitOps, setiap perubahan pada sistem dilakukan melalui \textit{pull request}, memberikan audit trail lengkap, kemudahan rollback, dan kolaborasi yang lebih baik \citep{gitops-principles}.

Penelitian ini berfokus pada implementasi GitOps untuk deployment aplikasi \textit{InvenioRDM} --- platform manajemen data riset \textit{open-source} --- pada klaster Kubernetes yang dikelola oleh Rancher. Studi kasus ini mencakup penggunaan ArgoCD sebagai \textit{GitOps controller}, Sealed Secrets untuk manajemen rahasia, dan Cloudflare Tunnel untuk akses eksternal yang aman.

\section{Rumusan Masalah}

Berdasarkan latar belakang di atas, rumusan masalah dalam penelitian ini adalah:

\begin{enumerate}
  \item Bagaimana merancang arsitektur GitOps untuk deployment aplikasi pada klaster Kubernetes?
  \item Bagaimana mengimplementasikan GitOps menggunakan ArgoCD untuk deployment otomatis aplikasi InvenioRDM?
  \item Bagaimana mengevaluasi efektivitas dan keamanan dari implementasi GitOps yang dibangun?
\end{enumerate}

\section{Batasan Masalah}

Agar penelitian ini tetap fokus dan terukur, batasan masalah yang ditetapkan adalah:

\begin{enumerate}
  \item Klaster Kubernetes yang digunakan dikelola oleh Rancher.
  \item \textit{GitOps controller} yang digunakan adalah ArgoCD.
  \item Aplikasi yang di-deploy adalah InvenioRDM dengan dependensinya (PostgreSQL, Redis, OpenSearch, MinIO).
  \item Evaluasi fokus pada aspek deployment, keamanan jaringan, dan manajemen rahasia.
\end{enumerate}

\section{Tujuan Penelitian}

Tujuan dari penelitian ini adalah:

\begin{enumerate}
  \item Merancang arsitektur GitOps untuk deployment aplikasi pada klaster Kubernetes.
  \item Mengimplementasikan GitOps menggunakan ArgoCD untuk deployment otomatis aplikasi InvenioRDM.
  \item Mengevaluasi efektivitas dan keamanan dari implementasi GitOps yang dibangun.
\end{enumerate}

\section{Manfaat Penelitian}

Penelitian ini diharapkan memberikan manfaat berikut:

\begin{enumerate}
  \item \textbf{Teoritis:} Menambah khazanah keilmuan mengenai implementasi GitOps pada Kubernetes, khususnya dalam konteks deployment aplikasi berbasis kontainer.
  \item \textbf{Praktis:} Memberikan panduan teknis dan \textit{blueprint} untuk implementasi GitOps pada klaster Kubernetes yang dapat diadopsi oleh organisasi atau peneliti lain.
\end{enumerate}

\section{Sistematika Penulisan}

Sistematika penulisan proposal tugas akhir ini adalah sebagai berikut:

\begin{itemize}
  \item \textbf{Bab I:} Pendahuluan --- berisi latar belakang, rumusan masalah, batasan masalah, tujuan, manfaat, dan sistematika penulisan.
  \item \textbf{Bab II:} Penelitian Terkait --- membahas state of the art dan hasil-hasil penelitian sebelumnya yang relevan.
  \item \textbf{Bab III:} Metode dan Rancangan Penelitian --- menjelaskan metodologi, perancangan sistem, bagan alir, dan tahapan implementasi.
  \item \textbf{Bab IV:} Jadwal Penelitian --- memuat Gantt Chart berbasis bulan untuk perencanaan waktu penelitian.
\end{itemize}
